\documentclass[a4paper,12pt]{article}

\usepackage[T1]{fontenc}
\usepackage[italian]{babel}
\usepackage{graphicx}
\usepackage{geometry}
\usepackage{tabularx} % Per l'ambiente tabularx (tabelle)
\usepackage{calc} % Sempre per le tabelle
\usepackage[table]{xcolor}
\usepackage[most]{tcolorbox}
\usepackage{fancyhdr}
\usepackage[hidelinks]{hyperref}
\usepackage{tikz}
\usepackage{pgf-pie}
\usepackage[utf8]{inputenc}
\usepackage{booktabs}
\usepackage{xcolor}  
\usepackage{float}
\usepackage{hyperref}
\usepackage{titlesec}
\usepackage{siunitx}
\usepackage{nicematrix}
\usepackage{colortbl}
\usepackage{eurosym}
\usepackage[dvipsnames]{xcolor}

\newcommand{\groupmail}{alittlebyte.swe@gmail.com}
\newcommand{\grouplogo}{../../../.github/site-src/assets/logo_alittlebyte.png}
\newcommand{\groupsymbol}{../../../.github/site-src/assets/solo_logo.png}

\geometry{
    left=2.5cm,
    right=2.5cm,
    top=2.5cm,
    bottom=2.5cm,
    headheight=331pt,
    includeheadfoot
}

\pagestyle{fancy}
\fancyhf{}

\renewcommand{\headrulewidth}{0pt}
\fancyhead{}

\renewcommand{\footrulewidth}{0.4pt}
\fancyfoot[L]{\includegraphics[height=0.6cm]{\groupsymbol}\hspace{0.45em}aLittleByte}
\fancyfoot[C]{\thepage}
\fancyfoot[R]{Verbale 2026-09-03}

\fancypagestyle{firstpage}{
    \fancyhf{}
    \renewcommand{\headrulewidth}{0pt}
    \renewcommand{\footrulewidth}{0.4pt}
    \fancyhead[C]{%
        \makebox[\textwidth][c]{%
            \begin{tabular}{c}
                \includegraphics[height=10cm]{\grouplogo}\\[1ex]
                \small\texttt{\groupmail}
            \end{tabular}%
        }
    }
    \fancyfoot[L]{\includegraphics[height=0.6cm]{\groupsymbol}\hspace{0.45em}aLittleByte}
    \fancyfoot[C]{\thepage}
    \fancyfoot[R]{Verbale 2026-09-03}
}

\providecommand{\glslink}[2]{\href{https://alittlebyte-19.github.io/Documentazione/glossario.html\#gls-#1}{\underline{#2}\textsuperscript{\scriptsize G}}}

\begin{document}

\thispagestyle{firstpage}

\vspace*{0.5cm}

\begin{center}
    {\LARGE \textbf{Verbale Esterno - 2026-09-03}}
\end{center}

\vspace{1.5cm}

\begin{center}
    \renewcommand{\arraystretch}{1.3}
    \begin{tcolorbox}[
        width=0.92\textwidth,
        colback=gray!6,
        colframe=gray!45,
        boxrule=0.5pt,
        arc=2mm,
        boxsep=0pt,
        left=8pt, right=8pt, top=8pt, bottom=8pt
    ]
        \begin{tabularx}{\linewidth}{>{\bfseries}p{3.5cm} X}
            Gruppo      & aLittleByte (gruppo 19) \\
            Redattore   & A. Oloni \\
            Verificatori& E. Bellet\\
            Destinatari & Prof. Tullio Vardanega, Prof. Riccardo Cardin \\
            Data        & 2026-09-03\\
        \end{tabularx}
    \end{tcolorbox}
\end{center}

\newpage

\newgeometry{left=2.5cm,right=2.5cm,top=2.5cm,bottom=2.5cm,includefoot}
\tableofcontents
\newpage

\section{Informazioni Generali}
\section*{Specifiche Documento}
\hrule
\vspace{1cm}

\begin{tabular}{>{\bfseries}p{5cm} | p{10cm}}
Luogo Riunione & Google Meet \\[6pt]

Orario (Inizio - Fine) & 15:15 -- 15.35 \\[6pt]

\glslink{redattore}{Redattore} &
A. Oloni\\[6pt]

\glslink{amministratore}{Amministratore} & 
L. Soligo\\[6pt]

\glslink{verificatore}{Verificatore} &
E. Bellet\\[6pt]

Partecipanti &
  A. Maule \\ &
  A. Oloni\\ &
  A. Gastaldello\\ &
  E. Bellet\\ & 
  D. Belluz\\ &
  L. Soligo\\ &
  V.Y. Kaneda Fernandes\\[6pt]

  Partecipanti Esterni &
  Luca Iuzzolino\\ &
  GianPaolo Ferrarin\\ [6pt]

\end{tabular}


\section{Ordine del Giorno}
    \begin{itemize}
    \item Presentazione dell'MVP (Minimum Viable Product) del progetto e dimostrazione live.
    \item Discussione e chiarimenti sulle metriche di affidabilità dell'applicativo.
    \item Valutazione e approvazione finale da parte dell'azienda referente.
    \end{itemize}


\section{Attività svolte}
\subsection{Presentazione dell'MVP del progetto e dimostrazione live.}
\textbf{Leonardo Soligo} ha guidato la presentazione dell'MVP rivolta ai referenti aziendali. Tramite la condivisione dello schermo, è stato mostrato l'applicativo in funzione. È stato confermato che lo stack tecnologico utilizzato è rimasto invariato rispetto alle scelte architetturali precedenti. Sono stati mostrati quindi:
\begin{itemize}
    \item Overview iniziale: panoramica su documenti da verificare, bozze da valutare e dati medi.
    \item Modulo AI Assistant: generazione automatica di testo, copertina e PDF finale, sistema di rating, storico documenti con filtri e metriche.
    \item Modulo Co-pilot: caricamento documenti, stato di validazione (automatica, da verificare, verificata manualmente) e metriche di utilizzo (es. documenti analizzati, tempi medi di elaborazione).
\end{itemize}

\subsection{Discussione e chiarimenti sulle metriche di affidabilità dell'applicativo.}
I referenti aziendali hanno richiesto chiarimenti specifici sul metodo di calcolo della percentuale di affidabilità/confidenza mostrata dal sistema. Leonardo Soligo ha quindi fornito le relative spiegazioni, illustrando anche visivamente le modalità con cui tali percentuali vengono determinate.

\subsection{Valutazione e approvazione finale da parte dell'azienda referente.}
L'azienda si è espressa in modo \textbf{positivo} riguardo al lavoro svolto e alle funzionalità presentate. Ha complimentato il team per i risultati raggiunti e ha \textbf{approvato} ufficialmente l'MVP presentato.
\section{To-Do List}
\begin{table}[h]
    \centering
    \begin{tabular}{|p{4.5cm}|p{5cm}|l|}
    \hline
        \rowcolor{lightgray}\textit{\textbf{Persona Incaricata}}  & \textbf{Task Assegnata} &\textbf{Link \glslink{issue}{Issue}} \\
    \hline
        \textit{Alex Oloni} & Redigere il verbale 2026-09-03 & \href{https://github.com/aLittleByte-19/Documentazione/issues/175}{\#175} \\
    \hline
    \end{tabular}
    \label{tab:placeholder}
\end{table}

\end{document}
